% ====================================================================
%  Demonstration of the functionality of the 'mul-beamer' class
%
%  This file is included by talk.tex and lecture.tex via
%    % ====================================================================
%  Demonstration of the functionality of the 'mul-beamer' class
%
%  This file is included by talk.tex and lecture.tex via
%    % ====================================================================
%  Demonstration of the functionality of the 'mul-beamer' class
%
%  This file is included by talk.tex and lecture.tex via
%    % ====================================================================
%  Demonstration of the functionality of the 'mul-beamer' class
%
%  This file is included by talk.tex and lecture.tex via
%    \input{functionality_demo}
%  Remove that line (and this file) when starting your own presentation.
% ====================================================================
%
% typesetting command names in the demo, e.g. \democmd{alert} -> \alert
\newcommand{\democmd}[1]{\texttt{\textbackslash #1}}
%
%
\section{Basics}
\subsection{Text and highlighting}
\begin{frame}{Text and highlighting}
	\begin{itemize}[<1->]% all items at once (see "Lists and overlays")
		\item \alert{important content} is highlighted with \democmd{alert\{\dots\}}
		\item \Alert{less important content} is highlighted with \democmd{Alert\{\dots\}}
		\item {\funneldisplay The MUL display font} (Funnel Display) is switched on with \democmd{funneldisplay}; it is used for the frame titles
		\item the text is set in Latin Modern Sans; class option \texttt{brandFont} switches to Host Grotesk, the MUL text font
		\item mathematics works as usual: $a^2 + b^2 = \Alert{c^2}$, \quad $\sin\alpha = \frac{b}{c}$
	\end{itemize}
\end{frame}
%
\subsection{Blocks and columns}
\begin{frame}{Blocks and columns}
	\begin{columns}[T]
		\begin{column}{0.48\textwidth}
			\begin{block}{A block}
				Standard beamer \texttt{block} in the MUL colours.
			\end{block}
			\begin{alertblock}{An alert block}
				For warnings and important statements.
			\end{alertblock}
		\end{column}
		\begin{column}{0.48\textwidth}
			\begin{exampleblock}{An example block}
				For examples.
			\end{exampleblock}
			\begin{block}{}
				A block without title.
			\end{block}
		\end{column}
	\end{columns}
\end{frame}
%
\subsection{Colours}
% colour swatch with its name, e.g. \democolour{digital_Turquoise}{white}
\newcommand{\democolour}[2]{%
	\tikz[baseline=-0.6ex]\node[fill=#1, text=#2, font=\tiny\ttfamily, minimum width=26mm, minimum height=5mm, inner sep=1pt] {\detokenize{#1}};}
\begin{frame}{Colours}
	The brand colours are defined with the names of the MUL style guide (also used by the poster template):
	\medskip
	\begin{columns}[T]
		\begin{column}{0.32\textwidth}
			\textbf{primary}\par\smallskip
			\democolour{digital_Turquoise}{white}\par
			\democolour{digital_GraphiteBlack}{white}\par
			\democolour{digital_LightGreen}{digital_GraphiteBlack}\par
			\democolour{digital_TechGreen}{digital_GraphiteBlack}\par
			\democolour{digital_White}{digital_GraphiteBlack}
		\end{column}
		\begin{column}{0.32\textwidth}
			\textbf{combination} (with a primary colour)\par\smallskip
			\democolour{digital_BTLight}{white}\par
			\democolour{digital_BTDark}{white}\par
			\democolour{digital_TGLight}{white}\par
			\democolour{digital_TGDark}{white}\par
			\democolour{digital_GWLight}{digital_GraphiteBlack}\par
			\democolour{digital_GWDark}{digital_GraphiteBlack}
		\end{column}
		\begin{column}{0.32\textwidth}
			\textbf{secondary} (study areas only)\par\smallskip
			\democolour{digital_Violet}{white}\par
			\democolour{digital_Orange}{white}\par
			\democolour{digital_Blue}{white}
		\end{column}
	\end{columns}
	\medskip
	Use them as usual, e.g.\ \texttt{\textbackslash textcolor\{digital\_Turquoise\}\{\dots\}}.
\end{frame}
%
%
\section{Lists and overlays}
\subsection{Lists}
\begin{frame}{Lists}
	By default, list items are uncovered one by one (\texttt{<+->} is the class default):
	\begin{itemize}
		\item first level items have squares
		\begin{itemize}
			\item lower levels have dashes
		\end{itemize}
		\example examples are marked with \democmd{example} instead of \democmd{item}
		\item last item
	\end{itemize}
	\bigskip
	Use \texttt{\textbackslash begin\{itemize\}[<1->]} to show all items at once:
	\begin{itemize}[<1->]
		\item all
		\item at
		\item once
	\end{itemize}
\end{frame}
%
\subsection{Overlays}
\begin{frame}{Overlays}
	\begin{itemize}[<1->]
		\item \democmd{only<2>\{\dots\}}: content exists on slide 2 only (takes no space otherwise)
		\item \democmd{uncover<3->\{\dots\}}: space for the content is always reserved, the content is shown from slide 3 on
	\end{itemize}
	\bigskip
	\only<2>{\Alert{This text is shown with \democmd{only<2>}.}}
	\uncover<3->{\alert{This text is shown with \democmd{uncover<3->}.}}
\end{frame}
%
%
\section{Figures}
\subsection{Figures and literature}
\literature{%
	A. Author, B. Author, and C. Author, {\em J. Example Sci.} {\bf 1}, 1 (2026)\\
	D. Author and E. Author, {\em Example Lett.} {\bf 2}, 100 (2025)}
\begin{frame}{Figures and literature}
	\begin{columns}[T]
		\begin{column}{0.55\textwidth}
			\centering
			\input{figs/demo_plot}
		\end{column}
		\begin{column}{0.4\textwidth}
			\centering
			\input{figs/demo_lattice}
		\end{column}
	\end{columns}
	\bigskip
	\begin{itemize}[<1->]
		\item figures are kept in the folder \texttt{figs}, graphics (logos, background) in \texttt{graphics}
		\item images are included with \texttt{\textbackslash includegraphics\{figs/\dots\}}; these demo figures are TikZ drawings included with \texttt{\textbackslash input}
		\item \democmd{literature\{\dots\}} (outside a frame) prints references above the footer of this and all following slides; clear it with \democmd{literature\{\}}
	\end{itemize}
\end{frame}
\literature{}
%
\subsection{Placing objects}
\begin{frame}{Placing objects with \texttt{overdraw}}
	\begin{itemize}[<1->]
		\item the \texttt{overdraw} environment is a TikZ picture with absolute coordinates: one unit is 1/16 of the slide width, the origin is the bottom-left corner of the slide (right of the sidebar, if shown)
		\item \democmd{framegrid} inside \texttt{overdraw} draws the coordinate grid for finding positions (slide 2) -- remove it in the final version
		\item slide 3: label and arrow placed with \texttt{overdraw}
	\end{itemize}
	\onslide<2>
	\begin{overdraw}
		\framegrid
	\end{overdraw}
	\onslide<3>
	\begin{overdraw}
		\node[draw=digital_Turquoise, very thick, rounded corners, fill=white, font=\small] (label) at (9,3) {placed at (9,3)};
		\draw[->, very thick, digital_Turquoise] (label) -- (5,4.5);
	\end{overdraw}
\end{frame}
%
%
\section{More tools}
\subsection{Hidden results}
\begin{frame}{Hidden results}
	\begin{itemize}[<1->]
		\item \democmd{showResults\{\dots\}} shows its content with the class option \texttt{showResults} (default), and leaves an empty space of the same size with \texttt{hideResults}
		\item useful for exercises in lecture handouts
	\end{itemize}
	\bigskip
	\begin{block}{Exercise}
		Solve $x^2 - 5x + 6 = 0$.\par\medskip
		Result: \showResults{$x_1 = 2$, \quad $x_2 = 3$}
	\end{block}
\end{frame}
%
\subsection{Frame logo}
% \setFrameLogo inside a group {...} changes the logo for the enclosed frames only;
% outside a group (or in the preamble) it changes it for all following frames
{%
\setFrameLogo{graphics/logo-dms.pdf}
\begin{frame}{Frame logo}
	\begin{itemize}[<1->]
		\item the logo in the top-right corner can be changed with \democmd{setFrameLogo\{graphics/\dots\}}
		\item here: the logo of the Department of Materials Science
	\end{itemize}
\end{frame}
}%
%
\subsection{Person slide}
\begin{frame}{Person slide}
	\begin{itemize}[<1->]
		\item \democmd{personframe[<photo>]\{<name>\}\{<role>\}\{<text>\}} (between frames) introduces a person, as in the \emph{Zur Person} layout of the MUL PowerPoint master
		\item the photo fills the left part of the slide (cropped as needed); without \texttt{<photo>}, a placeholder is shown (next slide)
		\item \texttt{<text>} can contain lists, e.g.\ a short CV
	\end{itemize}
\end{frame}
%
\personframe{Name Surname}{Function, Department of Materials Science}{%
	\begin{itemize}[<1->]
		\item since 2020: current position
		\item 2015--2020: previous position
		\item research interests: topic A, topic B
	\end{itemize}}
%
\subsection{Closing slide}
\begin{frame}{Closing slide}
	\begin{itemize}[<1->]
		\item \democmd{finalframe} (between frames) adds a closing slide in the dark style of the MUL slides: ``Thank you for your attention!'' (localised), the e-mail address set by \democmd{email}, and the claim
		\item \democmd{finalframe[\dots]} replaces the text, e.g.\ \texttt{\textbackslash finalframe[Questions?]}
		\item \democmd{setDarkBackground\{graphics/background\_dark.jpg\}} (in the preamble) uses the background of the official PowerPoint master instead of the CMS variant
		\item the class option \texttt{darkTitle} gives a title slide in the same dark style
	\end{itemize}
\end{frame}
%
\finalframe
%
% \lastslide: the total number of slides (option slideTotalNumber) counts only
% up to here, i.e., the following backup slides are not included
\lastslide
%
\section*{Backup}
\begin{frame}{Backup slide}
	\begin{itemize}[<1->]
		\item slides after \democmd{lastslide} are not counted in the total number of slides
		\item useful for backup slides
	\end{itemize}
\end{frame}

%  Remove that line (and this file) when starting your own presentation.
% ====================================================================
%
% typesetting command names in the demo, e.g. \democmd{alert} -> \alert
\newcommand{\democmd}[1]{\texttt{\textbackslash #1}}
%
%
\section{Basics}
\subsection{Text and highlighting}
\begin{frame}{Text and highlighting}
	\begin{itemize}[<1->]% all items at once (see "Lists and overlays")
		\item \alert{important content} is highlighted with \democmd{alert\{\dots\}}
		\item \Alert{less important content} is highlighted with \democmd{Alert\{\dots\}}
		\item {\funneldisplay The MUL display font} (Funnel Display) is switched on with \democmd{funneldisplay}; it is used for the frame titles
		\item the text is set in Latin Modern Sans; class option \texttt{brandFont} switches to Host Grotesk, the MUL text font
		\item mathematics works as usual: $a^2 + b^2 = \Alert{c^2}$, \quad $\sin\alpha = \frac{b}{c}$
	\end{itemize}
\end{frame}
%
\subsection{Blocks and columns}
\begin{frame}{Blocks and columns}
	\begin{columns}[T]
		\begin{column}{0.48\textwidth}
			\begin{block}{A block}
				Standard beamer \texttt{block} in the MUL colours.
			\end{block}
			\begin{alertblock}{An alert block}
				For warnings and important statements.
			\end{alertblock}
		\end{column}
		\begin{column}{0.48\textwidth}
			\begin{exampleblock}{An example block}
				For examples.
			\end{exampleblock}
			\begin{block}{}
				A block without title.
			\end{block}
		\end{column}
	\end{columns}
\end{frame}
%
\subsection{Colours}
% colour swatch with its name, e.g. \democolour{digital_Turquoise}{white}
\newcommand{\democolour}[2]{%
	\tikz[baseline=-0.6ex]\node[fill=#1, text=#2, font=\tiny\ttfamily, minimum width=26mm, minimum height=5mm, inner sep=1pt] {\detokenize{#1}};}
\begin{frame}{Colours}
	The brand colours are defined with the names of the MUL style guide (also used by the poster template):
	\medskip
	\begin{columns}[T]
		\begin{column}{0.32\textwidth}
			\textbf{primary}\par\smallskip
			\democolour{digital_Turquoise}{white}\par
			\democolour{digital_GraphiteBlack}{white}\par
			\democolour{digital_LightGreen}{digital_GraphiteBlack}\par
			\democolour{digital_TechGreen}{digital_GraphiteBlack}\par
			\democolour{digital_White}{digital_GraphiteBlack}
		\end{column}
		\begin{column}{0.32\textwidth}
			\textbf{combination} (with a primary colour)\par\smallskip
			\democolour{digital_BTLight}{white}\par
			\democolour{digital_BTDark}{white}\par
			\democolour{digital_TGLight}{white}\par
			\democolour{digital_TGDark}{white}\par
			\democolour{digital_GWLight}{digital_GraphiteBlack}\par
			\democolour{digital_GWDark}{digital_GraphiteBlack}
		\end{column}
		\begin{column}{0.32\textwidth}
			\textbf{secondary} (study areas only)\par\smallskip
			\democolour{digital_Violet}{white}\par
			\democolour{digital_Orange}{white}\par
			\democolour{digital_Blue}{white}
		\end{column}
	\end{columns}
	\medskip
	Use them as usual, e.g.\ \texttt{\textbackslash textcolor\{digital\_Turquoise\}\{\dots\}}.
\end{frame}
%
%
\section{Lists and overlays}
\subsection{Lists}
\begin{frame}{Lists}
	By default, list items are uncovered one by one (\texttt{<+->} is the class default):
	\begin{itemize}
		\item first level items have squares
		\begin{itemize}
			\item lower levels have dashes
		\end{itemize}
		\example examples are marked with \democmd{example} instead of \democmd{item}
		\item last item
	\end{itemize}
	\bigskip
	Use \texttt{\textbackslash begin\{itemize\}[<1->]} to show all items at once:
	\begin{itemize}[<1->]
		\item all
		\item at
		\item once
	\end{itemize}
\end{frame}
%
\subsection{Overlays}
\begin{frame}{Overlays}
	\begin{itemize}[<1->]
		\item \democmd{only<2>\{\dots\}}: content exists on slide 2 only (takes no space otherwise)
		\item \democmd{uncover<3->\{\dots\}}: space for the content is always reserved, the content is shown from slide 3 on
	\end{itemize}
	\bigskip
	\only<2>{\Alert{This text is shown with \democmd{only<2>}.}}
	\uncover<3->{\alert{This text is shown with \democmd{uncover<3->}.}}
\end{frame}
%
%
\section{Figures}
\subsection{Figures and literature}
\literature{%
	A. Author, B. Author, and C. Author, {\em J. Example Sci.} {\bf 1}, 1 (2026)\\
	D. Author and E. Author, {\em Example Lett.} {\bf 2}, 100 (2025)}
\begin{frame}{Figures and literature}
	\begin{columns}[T]
		\begin{column}{0.55\textwidth}
			\centering
			% demonstration graphics: a simple plot drawn with TikZ
% (uses the plot colours defined by mul-beamer.cls)
\begin{tikzpicture}[x=0.9cm, y=0.9cm]
	% axes
	\draw[->] (0,0) -- (6.4,0) node[below left] {\scriptsize time};
	\draw[->] (0,0) -- (0,3.4) node[above right] {\scriptsize signal};
	\foreach \x in {1,...,6} {\draw (\x,0) -- (\x,-0.08);}
	% data
	\draw[very thick, plotblue, domain=0:6, samples=80, smooth] plot (\x, {1.6 + 1.2*sin(1.5*\x r)*exp(-0.25*\x)});
	\draw[very thick, plotorange, dashed, domain=0:6] plot (\x, {0.4 + 0.35*\x});
	\foreach \x/\y in {0.5/0.7, 1.5/1.0, 2.5/1.3, 3.5/1.5, 4.5/2.1, 5.5/2.3} {\fill[plotred] (\x,\y) circle (2pt);}
	% legend
	\node[anchor=north east, font=\scriptsize, align=left] at (6.4,3.4) {%
		\textcolor{plotblue}{\rule[0.5ex]{4mm}{1pt}} model A\\
		\textcolor{plotorange}{\rule[0.5ex]{4mm}{1pt}} model B\\
		\textcolor{plotred}{$\bullet$} measurement};
\end{tikzpicture}%

		\end{column}
		\begin{column}{0.4\textwidth}
			\centering
			% demonstration graphics: a schematic two-component lattice drawn with TikZ
\begin{tikzpicture}[x=0.55cm, y=0.55cm]
	\draw[gray!50] (-0.5,-0.5) rectangle (5.5,4.5);
	\foreach \i in {0,...,5} {
		\foreach \j in {0,...,4} {
			\pgfmathtruncatemacro{\site}{mod(\i+\j,2)}
			\ifnum\site=0
				\shade[ball color=plotblue] (\i,\j) circle (0.3);
			\else
				\shade[ball color=plotred] (\i,\j) circle (0.22);
			\fi
		}
	}
	\node[below, font=\scriptsize] at (2.5,-0.5) {%
		\textcolor{plotblue}{$\bullet$} species A\quad
		\textcolor{plotred}{$\bullet$} species B};
\end{tikzpicture}%

		\end{column}
	\end{columns}
	\bigskip
	\begin{itemize}[<1->]
		\item figures are kept in the folder \texttt{figs}, graphics (logos, background) in \texttt{graphics}
		\item images are included with \texttt{\textbackslash includegraphics\{figs/\dots\}}; these demo figures are TikZ drawings included with \texttt{\textbackslash input}
		\item \democmd{literature\{\dots\}} (outside a frame) prints references above the footer of this and all following slides; clear it with \democmd{literature\{\}}
	\end{itemize}
\end{frame}
\literature{}
%
\subsection{Placing objects}
\begin{frame}{Placing objects with \texttt{overdraw}}
	\begin{itemize}[<1->]
		\item the \texttt{overdraw} environment is a TikZ picture with absolute coordinates: one unit is 1/16 of the slide width, the origin is the bottom-left corner of the slide (right of the sidebar, if shown)
		\item \democmd{framegrid} inside \texttt{overdraw} draws the coordinate grid for finding positions (slide 2) -- remove it in the final version
		\item slide 3: label and arrow placed with \texttt{overdraw}
	\end{itemize}
	\onslide<2>
	\begin{overdraw}
		\framegrid
	\end{overdraw}
	\onslide<3>
	\begin{overdraw}
		\node[draw=digital_Turquoise, very thick, rounded corners, fill=white, font=\small] (label) at (9,3) {placed at (9,3)};
		\draw[->, very thick, digital_Turquoise] (label) -- (5,4.5);
	\end{overdraw}
\end{frame}
%
%
\section{More tools}
\subsection{Hidden results}
\begin{frame}{Hidden results}
	\begin{itemize}[<1->]
		\item \democmd{showResults\{\dots\}} shows its content with the class option \texttt{showResults} (default), and leaves an empty space of the same size with \texttt{hideResults}
		\item useful for exercises in lecture handouts
	\end{itemize}
	\bigskip
	\begin{block}{Exercise}
		Solve $x^2 - 5x + 6 = 0$.\par\medskip
		Result: \showResults{$x_1 = 2$, \quad $x_2 = 3$}
	\end{block}
\end{frame}
%
\subsection{Frame logo}
% \setFrameLogo inside a group {...} changes the logo for the enclosed frames only;
% outside a group (or in the preamble) it changes it for all following frames
{%
\setFrameLogo{graphics/logo-dms.pdf}
\begin{frame}{Frame logo}
	\begin{itemize}[<1->]
		\item the logo in the top-right corner can be changed with \democmd{setFrameLogo\{graphics/\dots\}}
		\item here: the logo of the Department of Materials Science
	\end{itemize}
\end{frame}
}%
%
\subsection{Person slide}
\begin{frame}{Person slide}
	\begin{itemize}[<1->]
		\item \democmd{personframe[<photo>]\{<name>\}\{<role>\}\{<text>\}} (between frames) introduces a person, as in the \emph{Zur Person} layout of the MUL PowerPoint master
		\item the photo fills the left part of the slide (cropped as needed); without \texttt{<photo>}, a placeholder is shown (next slide)
		\item \texttt{<text>} can contain lists, e.g.\ a short CV
	\end{itemize}
\end{frame}
%
\personframe{Name Surname}{Function, Department of Materials Science}{%
	\begin{itemize}[<1->]
		\item since 2020: current position
		\item 2015--2020: previous position
		\item research interests: topic A, topic B
	\end{itemize}}
%
\subsection{Closing slide}
\begin{frame}{Closing slide}
	\begin{itemize}[<1->]
		\item \democmd{finalframe} (between frames) adds a closing slide in the dark style of the MUL slides: ``Thank you for your attention!'' (localised), the e-mail address set by \democmd{email}, and the claim
		\item \democmd{finalframe[\dots]} replaces the text, e.g.\ \texttt{\textbackslash finalframe[Questions?]}
		\item \democmd{setDarkBackground\{graphics/background\_dark.jpg\}} (in the preamble) uses the background of the official PowerPoint master instead of the CMS variant
		\item the class option \texttt{darkTitle} gives a title slide in the same dark style
	\end{itemize}
\end{frame}
%
\finalframe
%
% \lastslide: the total number of slides (option slideTotalNumber) counts only
% up to here, i.e., the following backup slides are not included
\lastslide
%
\section*{Backup}
\begin{frame}{Backup slide}
	\begin{itemize}[<1->]
		\item slides after \democmd{lastslide} are not counted in the total number of slides
		\item useful for backup slides
	\end{itemize}
\end{frame}

%  Remove that line (and this file) when starting your own presentation.
% ====================================================================
%
% typesetting command names in the demo, e.g. \democmd{alert} -> \alert
\newcommand{\democmd}[1]{\texttt{\textbackslash #1}}
%
%
\section{Basics}
\subsection{Text and highlighting}
\begin{frame}{Text and highlighting}
	\begin{itemize}[<1->]% all items at once (see "Lists and overlays")
		\item \alert{important content} is highlighted with \democmd{alert\{\dots\}}
		\item \Alert{less important content} is highlighted with \democmd{Alert\{\dots\}}
		\item {\funneldisplay The MUL display font} (Funnel Display) is switched on with \democmd{funneldisplay}; it is used for the frame titles
		\item the text is set in Latin Modern Sans; class option \texttt{brandFont} switches to Host Grotesk, the MUL text font
		\item mathematics works as usual: $a^2 + b^2 = \Alert{c^2}$, \quad $\sin\alpha = \frac{b}{c}$
	\end{itemize}
\end{frame}
%
\subsection{Blocks and columns}
\begin{frame}{Blocks and columns}
	\begin{columns}[T]
		\begin{column}{0.48\textwidth}
			\begin{block}{A block}
				Standard beamer \texttt{block} in the MUL colours.
			\end{block}
			\begin{alertblock}{An alert block}
				For warnings and important statements.
			\end{alertblock}
		\end{column}
		\begin{column}{0.48\textwidth}
			\begin{exampleblock}{An example block}
				For examples.
			\end{exampleblock}
			\begin{block}{}
				A block without title.
			\end{block}
		\end{column}
	\end{columns}
\end{frame}
%
\subsection{Colours}
% colour swatch with its name, e.g. \democolour{digital_Turquoise}{white}
\newcommand{\democolour}[2]{%
	\tikz[baseline=-0.6ex]\node[fill=#1, text=#2, font=\tiny\ttfamily, minimum width=26mm, minimum height=5mm, inner sep=1pt] {\detokenize{#1}};}
\begin{frame}{Colours}
	The brand colours are defined with the names of the MUL style guide (also used by the poster template):
	\medskip
	\begin{columns}[T]
		\begin{column}{0.32\textwidth}
			\textbf{primary}\par\smallskip
			\democolour{digital_Turquoise}{white}\par
			\democolour{digital_GraphiteBlack}{white}\par
			\democolour{digital_LightGreen}{digital_GraphiteBlack}\par
			\democolour{digital_TechGreen}{digital_GraphiteBlack}\par
			\democolour{digital_White}{digital_GraphiteBlack}
		\end{column}
		\begin{column}{0.32\textwidth}
			\textbf{combination} (with a primary colour)\par\smallskip
			\democolour{digital_BTLight}{white}\par
			\democolour{digital_BTDark}{white}\par
			\democolour{digital_TGLight}{white}\par
			\democolour{digital_TGDark}{white}\par
			\democolour{digital_GWLight}{digital_GraphiteBlack}\par
			\democolour{digital_GWDark}{digital_GraphiteBlack}
		\end{column}
		\begin{column}{0.32\textwidth}
			\textbf{secondary} (study areas only)\par\smallskip
			\democolour{digital_Violet}{white}\par
			\democolour{digital_Orange}{white}\par
			\democolour{digital_Blue}{white}
		\end{column}
	\end{columns}
	\medskip
	Use them as usual, e.g.\ \texttt{\textbackslash textcolor\{digital\_Turquoise\}\{\dots\}}.
\end{frame}
%
%
\section{Lists and overlays}
\subsection{Lists}
\begin{frame}{Lists}
	By default, list items are uncovered one by one (\texttt{<+->} is the class default):
	\begin{itemize}
		\item first level items have squares
		\begin{itemize}
			\item lower levels have dashes
		\end{itemize}
		\example examples are marked with \democmd{example} instead of \democmd{item}
		\item last item
	\end{itemize}
	\bigskip
	Use \texttt{\textbackslash begin\{itemize\}[<1->]} to show all items at once:
	\begin{itemize}[<1->]
		\item all
		\item at
		\item once
	\end{itemize}
\end{frame}
%
\subsection{Overlays}
\begin{frame}{Overlays}
	\begin{itemize}[<1->]
		\item \democmd{only<2>\{\dots\}}: content exists on slide 2 only (takes no space otherwise)
		\item \democmd{uncover<3->\{\dots\}}: space for the content is always reserved, the content is shown from slide 3 on
	\end{itemize}
	\bigskip
	\only<2>{\Alert{This text is shown with \democmd{only<2>}.}}
	\uncover<3->{\alert{This text is shown with \democmd{uncover<3->}.}}
\end{frame}
%
%
\section{Figures}
\subsection{Figures and literature}
\literature{%
	A. Author, B. Author, and C. Author, {\em J. Example Sci.} {\bf 1}, 1 (2026)\\
	D. Author and E. Author, {\em Example Lett.} {\bf 2}, 100 (2025)}
\begin{frame}{Figures and literature}
	\begin{columns}[T]
		\begin{column}{0.55\textwidth}
			\centering
			% demonstration graphics: a simple plot drawn with TikZ
% (uses the plot colours defined by mul-beamer.cls)
\begin{tikzpicture}[x=0.9cm, y=0.9cm]
	% axes
	\draw[->] (0,0) -- (6.4,0) node[below left] {\scriptsize time};
	\draw[->] (0,0) -- (0,3.4) node[above right] {\scriptsize signal};
	\foreach \x in {1,...,6} {\draw (\x,0) -- (\x,-0.08);}
	% data
	\draw[very thick, plotblue, domain=0:6, samples=80, smooth] plot (\x, {1.6 + 1.2*sin(1.5*\x r)*exp(-0.25*\x)});
	\draw[very thick, plotorange, dashed, domain=0:6] plot (\x, {0.4 + 0.35*\x});
	\foreach \x/\y in {0.5/0.7, 1.5/1.0, 2.5/1.3, 3.5/1.5, 4.5/2.1, 5.5/2.3} {\fill[plotred] (\x,\y) circle (2pt);}
	% legend
	\node[anchor=north east, font=\scriptsize, align=left] at (6.4,3.4) {%
		\textcolor{plotblue}{\rule[0.5ex]{4mm}{1pt}} model A\\
		\textcolor{plotorange}{\rule[0.5ex]{4mm}{1pt}} model B\\
		\textcolor{plotred}{$\bullet$} measurement};
\end{tikzpicture}%

		\end{column}
		\begin{column}{0.4\textwidth}
			\centering
			% demonstration graphics: a schematic two-component lattice drawn with TikZ
\begin{tikzpicture}[x=0.55cm, y=0.55cm]
	\draw[gray!50] (-0.5,-0.5) rectangle (5.5,4.5);
	\foreach \i in {0,...,5} {
		\foreach \j in {0,...,4} {
			\pgfmathtruncatemacro{\site}{mod(\i+\j,2)}
			\ifnum\site=0
				\shade[ball color=plotblue] (\i,\j) circle (0.3);
			\else
				\shade[ball color=plotred] (\i,\j) circle (0.22);
			\fi
		}
	}
	\node[below, font=\scriptsize] at (2.5,-0.5) {%
		\textcolor{plotblue}{$\bullet$} species A\quad
		\textcolor{plotred}{$\bullet$} species B};
\end{tikzpicture}%

		\end{column}
	\end{columns}
	\bigskip
	\begin{itemize}[<1->]
		\item figures are kept in the folder \texttt{figs}, graphics (logos, background) in \texttt{graphics}
		\item images are included with \texttt{\textbackslash includegraphics\{figs/\dots\}}; these demo figures are TikZ drawings included with \texttt{\textbackslash input}
		\item \democmd{literature\{\dots\}} (outside a frame) prints references above the footer of this and all following slides; clear it with \democmd{literature\{\}}
	\end{itemize}
\end{frame}
\literature{}
%
\subsection{Placing objects}
\begin{frame}{Placing objects with \texttt{overdraw}}
	\begin{itemize}[<1->]
		\item the \texttt{overdraw} environment is a TikZ picture with absolute coordinates: one unit is 1/16 of the slide width, the origin is the bottom-left corner of the slide (right of the sidebar, if shown)
		\item \democmd{framegrid} inside \texttt{overdraw} draws the coordinate grid for finding positions (slide 2) -- remove it in the final version
		\item slide 3: label and arrow placed with \texttt{overdraw}
	\end{itemize}
	\onslide<2>
	\begin{overdraw}
		\framegrid
	\end{overdraw}
	\onslide<3>
	\begin{overdraw}
		\node[draw=digital_Turquoise, very thick, rounded corners, fill=white, font=\small] (label) at (9,3) {placed at (9,3)};
		\draw[->, very thick, digital_Turquoise] (label) -- (5,4.5);
	\end{overdraw}
\end{frame}
%
%
\section{More tools}
\subsection{Hidden results}
\begin{frame}{Hidden results}
	\begin{itemize}[<1->]
		\item \democmd{showResults\{\dots\}} shows its content with the class option \texttt{showResults} (default), and leaves an empty space of the same size with \texttt{hideResults}
		\item useful for exercises in lecture handouts
	\end{itemize}
	\bigskip
	\begin{block}{Exercise}
		Solve $x^2 - 5x + 6 = 0$.\par\medskip
		Result: \showResults{$x_1 = 2$, \quad $x_2 = 3$}
	\end{block}
\end{frame}
%
\subsection{Frame logo}
% \setFrameLogo inside a group {...} changes the logo for the enclosed frames only;
% outside a group (or in the preamble) it changes it for all following frames
{%
\setFrameLogo{graphics/logo-dms.pdf}
\begin{frame}{Frame logo}
	\begin{itemize}[<1->]
		\item the logo in the top-right corner can be changed with \democmd{setFrameLogo\{graphics/\dots\}}
		\item here: the logo of the Department of Materials Science
	\end{itemize}
\end{frame}
}%
%
\subsection{Person slide}
\begin{frame}{Person slide}
	\begin{itemize}[<1->]
		\item \democmd{personframe[<photo>]\{<name>\}\{<role>\}\{<text>\}} (between frames) introduces a person, as in the \emph{Zur Person} layout of the MUL PowerPoint master
		\item the photo fills the left part of the slide (cropped as needed); without \texttt{<photo>}, a placeholder is shown (next slide)
		\item \texttt{<text>} can contain lists, e.g.\ a short CV
	\end{itemize}
\end{frame}
%
\personframe{Name Surname}{Function, Department of Materials Science}{%
	\begin{itemize}[<1->]
		\item since 2020: current position
		\item 2015--2020: previous position
		\item research interests: topic A, topic B
	\end{itemize}}
%
\subsection{Closing slide}
\begin{frame}{Closing slide}
	\begin{itemize}[<1->]
		\item \democmd{finalframe} (between frames) adds a closing slide in the dark style of the MUL slides: ``Thank you for your attention!'' (localised), the e-mail address set by \democmd{email}, and the claim
		\item \democmd{finalframe[\dots]} replaces the text, e.g.\ \texttt{\textbackslash finalframe[Questions?]}
		\item \democmd{setDarkBackground\{graphics/background\_dark.jpg\}} (in the preamble) uses the background of the official PowerPoint master instead of the CMS variant
		\item the class option \texttt{darkTitle} gives a title slide in the same dark style
	\end{itemize}
\end{frame}
%
\finalframe
%
% \lastslide: the total number of slides (option slideTotalNumber) counts only
% up to here, i.e., the following backup slides are not included
\lastslide
%
\section*{Backup}
\begin{frame}{Backup slide}
	\begin{itemize}[<1->]
		\item slides after \democmd{lastslide} are not counted in the total number of slides
		\item useful for backup slides
	\end{itemize}
\end{frame}

%  Remove that line (and this file) when starting your own presentation.
% ====================================================================
%
% typesetting command names in the demo, e.g. \democmd{alert} -> \alert
\newcommand{\democmd}[1]{\texttt{\textbackslash #1}}
%
%
\section{Basics}
\subsection{Text and highlighting}
\begin{frame}{Text and highlighting}
	\begin{itemize}[<1->]% all items at once (see "Lists and overlays")
		\item \alert{important content} is highlighted with \democmd{alert\{\dots\}}
		\item \Alert{less important content} is highlighted with \democmd{Alert\{\dots\}}
		\item {\funneldisplay The MUL display font} (Funnel Display) is switched on with \democmd{funneldisplay}; it is used for the frame titles
		\item the text is set in Latin Modern Sans; class option \texttt{brandFont} switches to Host Grotesk, the MUL text font
		\item mathematics works as usual: $a^2 + b^2 = \Alert{c^2}$, \quad $\sin\alpha = \frac{b}{c}$
	\end{itemize}
\end{frame}
%
\subsection{Blocks and columns}
\begin{frame}{Blocks and columns}
	\begin{columns}[T]
		\begin{column}{0.48\textwidth}
			\begin{block}{A block}
				Standard beamer \texttt{block} in the MUL colours.
			\end{block}
			\begin{alertblock}{An alert block}
				For warnings and important statements.
			\end{alertblock}
		\end{column}
		\begin{column}{0.48\textwidth}
			\begin{exampleblock}{An example block}
				For examples.
			\end{exampleblock}
			\begin{block}{}
				A block without title.
			\end{block}
		\end{column}
	\end{columns}
\end{frame}
%
\subsection{Colours}
% colour swatch with its name, e.g. \democolour{digital_Turquoise}{white}
\newcommand{\democolour}[2]{%
	\tikz[baseline=-0.6ex]\node[fill=#1, text=#2, font=\tiny\ttfamily, minimum width=26mm, minimum height=5mm, inner sep=1pt] {\detokenize{#1}};}
\begin{frame}{Colours}
	The brand colours are defined with the names of the MUL style guide (also used by the poster template):
	\medskip
	\begin{columns}[T]
		\begin{column}{0.32\textwidth}
			\textbf{primary}\par\smallskip
			\democolour{digital_Turquoise}{white}\par
			\democolour{digital_GraphiteBlack}{white}\par
			\democolour{digital_LightGreen}{digital_GraphiteBlack}\par
			\democolour{digital_TechGreen}{digital_GraphiteBlack}\par
			\democolour{digital_White}{digital_GraphiteBlack}
		\end{column}
		\begin{column}{0.32\textwidth}
			\textbf{combination} (with a primary colour)\par\smallskip
			\democolour{digital_BTLight}{white}\par
			\democolour{digital_BTDark}{white}\par
			\democolour{digital_TGLight}{white}\par
			\democolour{digital_TGDark}{white}\par
			\democolour{digital_GWLight}{digital_GraphiteBlack}\par
			\democolour{digital_GWDark}{digital_GraphiteBlack}
		\end{column}
		\begin{column}{0.32\textwidth}
			\textbf{secondary} (study areas only)\par\smallskip
			\democolour{digital_Violet}{white}\par
			\democolour{digital_Orange}{white}\par
			\democolour{digital_Blue}{white}
		\end{column}
	\end{columns}
	\medskip
	Use them as usual, e.g.\ \texttt{\textbackslash textcolor\{digital\_Turquoise\}\{\dots\}}.
\end{frame}
%
%
\section{Lists and overlays}
\subsection{Lists}
\begin{frame}{Lists}
	By default, list items are uncovered one by one (\texttt{<+->} is the class default):
	\begin{itemize}
		\item first level items have squares
		\begin{itemize}
			\item lower levels have dashes
		\end{itemize}
		\example examples are marked with \democmd{example} instead of \democmd{item}
		\item last item
	\end{itemize}
	\bigskip
	Use \texttt{\textbackslash begin\{itemize\}[<1->]} to show all items at once:
	\begin{itemize}[<1->]
		\item all
		\item at
		\item once
	\end{itemize}
\end{frame}
%
\subsection{Overlays}
\begin{frame}{Overlays}
	\begin{itemize}[<1->]
		\item \democmd{only<2>\{\dots\}}: content exists on slide 2 only (takes no space otherwise)
		\item \democmd{uncover<3->\{\dots\}}: space for the content is always reserved, the content is shown from slide 3 on
	\end{itemize}
	\bigskip
	\only<2>{\Alert{This text is shown with \democmd{only<2>}.}}
	\uncover<3->{\alert{This text is shown with \democmd{uncover<3->}.}}
\end{frame}
%
%
\section{Figures}
\subsection{Figures and literature}
\literature{%
	A. Author, B. Author, and C. Author, {\em J. Example Sci.} {\bf 1}, 1 (2026)\\
	D. Author and E. Author, {\em Example Lett.} {\bf 2}, 100 (2025)}
\begin{frame}{Figures and literature}
	\begin{columns}[T]
		\begin{column}{0.55\textwidth}
			\centering
			% demonstration graphics: a simple plot drawn with TikZ
% (uses the plot colours defined by mul-beamer.cls)
\begin{tikzpicture}[x=0.9cm, y=0.9cm]
	% axes
	\draw[->] (0,0) -- (6.4,0) node[below left] {\scriptsize time};
	\draw[->] (0,0) -- (0,3.4) node[above right] {\scriptsize signal};
	\foreach \x in {1,...,6} {\draw (\x,0) -- (\x,-0.08);}
	% data
	\draw[very thick, plotblue, domain=0:6, samples=80, smooth] plot (\x, {1.6 + 1.2*sin(1.5*\x r)*exp(-0.25*\x)});
	\draw[very thick, plotorange, dashed, domain=0:6] plot (\x, {0.4 + 0.35*\x});
	\foreach \x/\y in {0.5/0.7, 1.5/1.0, 2.5/1.3, 3.5/1.5, 4.5/2.1, 5.5/2.3} {\fill[plotred] (\x,\y) circle (2pt);}
	% legend
	\node[anchor=north east, font=\scriptsize, align=left] at (6.4,3.4) {%
		\textcolor{plotblue}{\rule[0.5ex]{4mm}{1pt}} model A\\
		\textcolor{plotorange}{\rule[0.5ex]{4mm}{1pt}} model B\\
		\textcolor{plotred}{$\bullet$} measurement};
\end{tikzpicture}%

		\end{column}
		\begin{column}{0.4\textwidth}
			\centering
			% demonstration graphics: a schematic two-component lattice drawn with TikZ
\begin{tikzpicture}[x=0.55cm, y=0.55cm]
	\draw[gray!50] (-0.5,-0.5) rectangle (5.5,4.5);
	\foreach \i in {0,...,5} {
		\foreach \j in {0,...,4} {
			\pgfmathtruncatemacro{\site}{mod(\i+\j,2)}
			\ifnum\site=0
				\shade[ball color=plotblue] (\i,\j) circle (0.3);
			\else
				\shade[ball color=plotred] (\i,\j) circle (0.22);
			\fi
		}
	}
	\node[below, font=\scriptsize] at (2.5,-0.5) {%
		\textcolor{plotblue}{$\bullet$} species A\quad
		\textcolor{plotred}{$\bullet$} species B};
\end{tikzpicture}%

		\end{column}
	\end{columns}
	\bigskip
	\begin{itemize}[<1->]
		\item figures are kept in the folder \texttt{figs}, graphics (logos, background) in \texttt{graphics}
		\item images are included with \texttt{\textbackslash includegraphics\{figs/\dots\}}; these demo figures are TikZ drawings included with \texttt{\textbackslash input}
		\item \democmd{literature\{\dots\}} (outside a frame) prints references above the footer of this and all following slides; clear it with \democmd{literature\{\}}
	\end{itemize}
\end{frame}
\literature{}
%
\subsection{Placing objects}
\begin{frame}{Placing objects with \texttt{overdraw}}
	\begin{itemize}[<1->]
		\item the \texttt{overdraw} environment is a TikZ picture with absolute coordinates: one unit is 1/16 of the slide width, the origin is the bottom-left corner of the slide (right of the sidebar, if shown)
		\item \democmd{framegrid} inside \texttt{overdraw} draws the coordinate grid for finding positions (slide 2) -- remove it in the final version
		\item slide 3: label and arrow placed with \texttt{overdraw}
	\end{itemize}
	\onslide<2>
	\begin{overdraw}
		\framegrid
	\end{overdraw}
	\onslide<3>
	\begin{overdraw}
		\node[draw=digital_Turquoise, very thick, rounded corners, fill=white, font=\small] (label) at (9,3) {placed at (9,3)};
		\draw[->, very thick, digital_Turquoise] (label) -- (5,4.5);
	\end{overdraw}
\end{frame}
%
%
\section{More tools}
\subsection{Hidden results}
\begin{frame}{Hidden results}
	\begin{itemize}[<1->]
		\item \democmd{showResults\{\dots\}} shows its content with the class option \texttt{showResults} (default), and leaves an empty space of the same size with \texttt{hideResults}
		\item useful for exercises in lecture handouts
	\end{itemize}
	\bigskip
	\begin{block}{Exercise}
		Solve $x^2 - 5x + 6 = 0$.\par\medskip
		Result: \showResults{$x_1 = 2$, \quad $x_2 = 3$}
	\end{block}
\end{frame}
%
\subsection{Frame logo}
% \setFrameLogo inside a group {...} changes the logo for the enclosed frames only;
% outside a group (or in the preamble) it changes it for all following frames
{%
\setFrameLogo{graphics/logo-dms.pdf}
\begin{frame}{Frame logo}
	\begin{itemize}[<1->]
		\item the logo in the top-right corner can be changed with \democmd{setFrameLogo\{graphics/\dots\}}
		\item here: the logo of the Department of Materials Science
	\end{itemize}
\end{frame}
}%
%
\subsection{Person slide}
\begin{frame}{Person slide}
	\begin{itemize}[<1->]
		\item \democmd{personframe[<photo>]\{<name>\}\{<role>\}\{<text>\}} (between frames) introduces a person, as in the \emph{Zur Person} layout of the MUL PowerPoint master
		\item the photo fills the left part of the slide (cropped as needed); without \texttt{<photo>}, a placeholder is shown (next slide)
		\item \texttt{<text>} can contain lists, e.g.\ a short CV
	\end{itemize}
\end{frame}
%
\personframe{Name Surname}{Function, Department of Materials Science}{%
	\begin{itemize}[<1->]
		\item since 2020: current position
		\item 2015--2020: previous position
		\item research interests: topic A, topic B
	\end{itemize}}
%
\subsection{Closing slide}
\begin{frame}{Closing slide}
	\begin{itemize}[<1->]
		\item \democmd{finalframe} (between frames) adds a closing slide in the dark style of the MUL slides: ``Thank you for your attention!'' (localised), the e-mail address set by \democmd{email}, and the claim
		\item \democmd{finalframe[\dots]} replaces the text, e.g.\ \texttt{\textbackslash finalframe[Questions?]}
		\item \democmd{setDarkBackground\{graphics/background\_dark.jpg\}} (in the preamble) uses the background of the official PowerPoint master instead of the CMS variant
		\item the class option \texttt{darkTitle} gives a title slide in the same dark style
	\end{itemize}
\end{frame}
%
\finalframe
%
% \lastslide: the total number of slides (option slideTotalNumber) counts only
% up to here, i.e., the following backup slides are not included
\lastslide
%
\section*{Backup}
\begin{frame}{Backup slide}
	\begin{itemize}[<1->]
		\item slides after \democmd{lastslide} are not counted in the total number of slides
		\item useful for backup slides
	\end{itemize}
\end{frame}
